%% Submitted to: 
%https://www.sciencedirect.com/journal/results-in-engineering

\documentclass[preprint,12pt]{elsarticle}

\usepackage[numbers]{natbib}

%%%Author macros
\def\tsc#1{\csdef{#1}{\textsc{\lowercase{#1}}\xspace}}
\tsc{WGM}
\tsc{QE}
\tsc{EP}
\tsc{PMS}
\tsc{BEC}
\tsc{DE}
%%%
\usepackage{caption}

\usepackage{float}
\usepackage{amssymb,amsmath}
\usepackage{url}
\usepackage{xcolor}
\usepackage{bm}
\usepackage{placeins}
\usepackage{hyperref}
\usepackage{longtable} 
\usepackage[capitalise]{cleveref}
\usepackage{svg}
\usepackage{pifont}
\usepackage{lineno}
\usepackage{placeins}

\newcommand{\cmark}{\ding{51}}
\newcommand{\xmark}{\ding{55}}

\usepackage{array}
\newcolumntype{C}[1]{>{\centering\let\newline\\\arraybackslash\hspace{0pt}}m{#1}}


\usepackage{tikz}
\usetikzlibrary{shapes.geometric, arrows.meta, positioning, calc}

\definecolor{inputgreen}{RGB}{230, 245, 230} % Green for definitions
\definecolor{processblue}{RGB}{230, 245, 255} % Blue for physics/simulation
\definecolor{outputorange}{RGB}{255, 245, 230} % Orange for final data

%%Defining some useful variables 
\providecommand{\ve}[1]{{\bm{#1}}}
\providecommand{\mat}[1]{{\bm{#1}}} 
\providecommand{\var}[1]

\DeclareMathOperator{\subconj}{\negthinspace\subset\negthinspace }
\DeclareMathOperator{\en}{\!\,\in\!\,}
\DeclareMathOperator{\igual}{\!\,=\!\,}
\DeclareMathOperator{\dist}{\operatorname{d}}
\DeclareMathOperator{\xx}{\negthickspace\times\negthickspace}
\providecommand{\s}[1]{\negthinspace#1\negthinspace}

\begin{document}

\newcommand{\Real}{\mathbb{R}}
\newcommand{\N}{\mathbb{N}}
\newcommand{\Z}{\mathbb{Z}}
\newcommand{\X}{\mathcal{X}}
\newcommand{\Y}{\mathcal{Y}}
\newcommand{\boldV}{\mathbf{V}}
\newcommand{\dataset}{{\cal D}}
\newcommand{\gauss}{\mathcal{N}} 

\providecommand{\ch}[1]{\textcolor{blue!50}{#1}}

\let\WriteBookmarks\relax
\def\floatpagepagefraction{1}
\def\textpagefraction{.001}

%% Journal name (optional)
\journal{Results in Engineering}

\begin{frontmatter}

\begin{highlights}
\item We employ an atomistic Hamiltonian to generate thermodynamically relaxed 2D magnetic domain images.
\item A regression-based deep learning framework accurately predicts Hamiltonian parameters from images.
\item A novel metric, Regression Activation Maps, is introduced to spatially explain continuous prediction errors.
\item The framework isolates critical degenerate textures, bridging black-box AI with physical laws.
\end{highlights}

\end{frontmatter}


\end{document}